\documentclass{situationist-preprint}

\title{The Amoralist Challenge}
\msid{MTE-07-b}
\msfamily{MTE-07}
\msdiscipline{Metaethics}
\mstype{Companion manuscript}
\mskeywords{motivational internalism; amoralist; practical rationality; externalism; akrasia}

\begin{document}
\maketitle

\begin{abstract}
Externalists argue from the amoralist, who sincerely judges that she ought to act and is wholly unmoved, that moral judgment does not necessarily motivate. MTE-07-a grants that the amoralist is possible and argues that she is practically irrational. The strongest form of the challenge describes an amoralist who is rational in every other respect, so that the charge of irrationality looks ad hoc. We argue that practical rationality is assessed belief by belief, not globally, and that the amoralist's failure is of the same type as the akratic's, differing only in degree. The challenge therefore leaves the conditional internalism of MTE-07 intact.
\end{abstract}

\section{The challenge and the reply}\label{s:ch}

\citet{Brink1986} argues that the amoralist is conceivable and that internalism must therefore be false; \citet{Svavarsdottir1999} sharpens the case with an agent who makes moral judgments competently, is unmoved by them, and is otherwise an ordinary, reasonable person. MTE-07-a accepts both points about possibility: its Proposition~3.2 states that amoralists are possible. What it denies is that the amoralist is practically rational (Corollary~3.3 of MTE-07-a). Its thesis is that a practically rational agent who judges that she morally ought to act is motivated to some degree (Theorem~3.1 of MTE-07-a), where practical rationality is the disposition of motivation to respond to one's beliefs about one's reasons (Definition~2.3 of MTE-07-a).

The challenge now becomes: Svavarsdóttir's agent is rational in every other respect, so to call her irrational is to define rationality so as to save the thesis. The reply is that Definition~2.3 of MTE-07-a is a condition on each belief about reasons, not on the agent as a whole. An agent can satisfy it with respect to her beliefs about her prudential, epistemic and aesthetic reasons and fail it with respect to one class of beliefs. Her failure is then local, as an agent who reasons well except about her own health is locally irrational. The amoralist's general reasonableness is no evidence of rationality in the respect at issue.

The comparison with weakness of will makes this precise. The akratic agent judges that she ought to act, is motivated by her judgment, but not enough, and acts otherwise, as ACT-04 allows. The amoralist judges that she ought to act and is not motivated at all. Both fail to be moved as their own judgments call for; the amoralist fails more completely. If the akratic's failure is a failure of practical rationality, as is generally granted, the amoralist's is too.

\begin{objection}
The amoralist does not believe that there are moral reasons for her; she believes only that the act is what morality requires. Her judgment is in inverted commas \citep[ch.~3]{Smith1994}.
\end{objection}

\begin{reply}
Then she is not the amoralist of the challenge. An agent who judges only that others call the act required does not judge that she ought to do it, and her lack of motivation is consistent with internalism of every form. The challenge needs an agent whose judgment is sincere, and by Principle P2 of MTE-07-a, to judge sincerely that one morally ought to act is to judge that there are decisive moral reasons for one to act. That belief is what her motivation fails to answer.
\end{reply}

\section{Conclusion}

The amoralist is possible, and she may be reasonable in most of her life; but in the respect at issue her motivations do not respond to her own beliefs about her reasons. The challenge refutes unconditional internalism and leaves standing the thesis that moral judgment necessarily motivates rational agents.

\begin{thebibliography}{99}\small
\bibitem[Brink(1986)]{Brink1986} Brink, D.~O. (1986). Externalist moral realism. \emph{The Southern Journal of Philosophy}, 24(S1), 23--41.
\bibitem[Smith(1994)]{Smith1994} Smith, M. (1994). \emph{The Moral Problem}. Blackwell.
\bibitem[Svavarsd\'ottir(1999)]{Svavarsdottir1999} Svavarsd\'ottir, S. (1999). Moral cognitivism and motivation. \emph{The Philosophical Review}, 108(2), 161--219.
\end{thebibliography}

\end{document}
